% ====================================================================
% introduction.tex (dissertation example)
% ====================================================================

\section{Introduction}
\label{sec:introduction}

\subsection{Motivation}

\lettrine[lines=2, slope=0em, loversize=0.14, nindent=0pt, lhang=0.3, lraise=0]{M}{OST}
low-cost microcontrollers include an analogue-to-digital converter (ADC) with
a resolution of 12 bits. With a \qty{3.3}{\volt} reference this gives a step
size of \qty{0.8}{\milli\volt}, which is adequate for many measurements, but
not for sensors whose output moves by only a small part of the input range. A
load cell or a thermocouple, for example, may need steps of tens of microvolts
to give a useful reading.

The usual answer is to add an external converter with a higher resolution,
which adds cost, board area and another device to configure. An alternative is
to take many readings with the converter already present and average them.
This is known as oversampling, and in principle it gives one extra bit of
resolution for every factor of four in the number of readings
averaged~\cite{atmel_avr121_2005}. It costs nothing in hardware, but the
output rate falls by the same factor, so it suits signals that change slowly.

The method has a condition attached, which is that the error in each reading
must differ from the last. If every reading of a steady input returns the same
code, their average is that same code and nothing has been gained. Some noise
at the input is therefore needed, and application notes generally say that it
should be around one LSB, without showing what happens when there is less.
That gap is what this project sets out to fill.

\subsection{Aim and objectives}

The aim of the project is to establish how much input noise an ideal 12-bit
converter needs for oversampling and averaging to deliver the resolution that
theory predicts. The objectives are:

\begin{enumerate}
  \item To review the theory of quantisation noise and of oversampling, and
        to state the predicted resolution as an equation that can be tested.
  \item To build a simulation model of an ideal converter followed by an
        averaging stage, and to check it against a case with a known answer.
  \item To measure the effective resolution for oversampling ratios from 1
        to 256, for a sine wave and for static inputs, at several noise
        levels.
  \item To compare the results with the theory, and to state the conditions
        under which the two agree.
\end{enumerate}

\subsection{Structure of the dissertation}

\Cref{sec:background} reviews the theory and the published guidance.
\Cref{sec:method} describes the simulation model and the two tests carried
out with it. The results are presented in \cref{sec:results} and discussed in
\cref{sec:discussion}, along with the limitations of the model.
\Cref{sec:conclusion} gives the conclusions and suggests further work. The
source code of the simulation is listed in Appendix~\ref{app:source}, and the
full table of results for the static test is in Appendix~\ref{app:results}.
